% ---- pg-388 （书页 376） ----

\[
=\mathbf{L}\mathbf{F}(\mathbf{L}\mathbf{F})'+\boldsymbol{\varepsilon}(\mathbf{L}\mathbf{F})'+\mathbf{L}\mathbf{F}\boldsymbol{\varepsilon}'+\boldsymbol{\varepsilon}\boldsymbol{\varepsilon}'
\]
故按照\hyperref[eq:9-3]{式(9-3)}有
\[
\begin{aligned}
\boldsymbol{\Sigma}&=\mathrm{Cov}(\mathbf{X})=E(\mathbf{X}-\boldsymbol{\mu})(\mathbf{X}-\boldsymbol{\mu})'\\
&=\mathbf{L}E(\mathbf{F}\mathbf{F}')\mathbf{L}'+E(\boldsymbol{\varepsilon}\mathbf{F}')\mathbf{L}'+\mathbf{L}E(\mathbf{F}\boldsymbol{\varepsilon}')+E(\boldsymbol{\varepsilon}\boldsymbol{\varepsilon}')\\
&=\mathbf{L}\mathbf{L}'+\boldsymbol{\Psi}
\end{aligned}
\]
而且, $\mathrm{Cov}(\boldsymbol{\varepsilon},\mathbf{F})=E(\boldsymbol{\varepsilon}\mathbf{F}')=\mathbf{0}$

由模型(9-4), 也有 $(\mathbf{X}-\boldsymbol{\mu})\mathbf{F}'=(\mathbf{L}\mathbf{F}+\boldsymbol{\varepsilon})\mathbf{F}'=\mathbf{L}\mathbf{F}\mathbf{F}'+\boldsymbol{\varepsilon}\mathbf{F}'$, 故 $\mathrm{Cov}(\mathbf{X},\mathbf{F})=E(\mathbf{X}-\boldsymbol{\mu})\mathbf{F}'=\mathbf{L}E(\mathbf{F}\mathbf{F}')+E(\boldsymbol{\varepsilon}\mathbf{F}')=\mathbf{L}$.

\begin{mvbox}
\begin{center}
\textbf{正交因子模型的协方差结构}
\end{center}
1. $\mathrm{Cov}(\mathbf{X})=\mathbf{L}\mathbf{L}'+\boldsymbol{\Psi}$

或
\phantomsection\label{eq:9-5}
\begin{equation}
\begin{aligned}
\mathrm{Var}(X_i)&=l_{i1}^2+\cdots+l_{im}^2+\psi_i\\
\mathrm{Cov}(X_i,X_k)&=l_{i1}l_{k1}+\cdots+l_{im}l_{km}
\end{aligned}
\tag{9-5}
\end{equation}
2. $\mathrm{Cov}(\mathbf{X},\mathbf{F})=\mathbf{L}$

或
\[
\mathrm{Cov}(X_i,F_j)=l_{ij}
\]
\end{mvbox}

模型 $\mathbf{X}-\boldsymbol{\mu}=\mathbf{L}\mathbf{F}+\boldsymbol{\varepsilon}$ 在公共因子中是线性的. 如果 $p$ 个响应 $\mathbf{X}$ 事实上是与潜在因子有关, 但这关系为非线性的, 比如为 $X_1-\mu_1=l_{11}F_1F_3+\varepsilon_1$, $X_2-\mu_2=l_{21}F_2F_3+\varepsilon_2$, 那么仅由\hyperref[eq:9-5]{式(9-5)}给出的协方差结构 $\mathbf{L}\mathbf{L}'+\boldsymbol{\Psi}$ 可能还是不够的. 线性假设十分重要, 它是传统因子模型公式表达中固有的假设.

由 $m$ 个公共因子贡献的第 $i$ 个变量的方差部分, 叫做第 $i$ 个\mvbold{共性方差}. 属于特殊因子的 $\mathrm{Var}(X_i)=\sigma_{ii}$ 部分, 常称为\mvbold{独特方差}或\mvbold{特殊方差}. 用 $h_i^2$ 记第 $i$ 个共性方差, 从\hyperref[eq:9-5]{式(9-5)}我们看到
\[
\underbrace{\sigma_{ii}}_{\mathrm{Var}(X_i)}=\underbrace{l_{i1}^2+l_{i2}^2+\cdots+l_{im}^2}_{\text{共性方差}}+\underbrace{\psi_i}_{\text{特殊方差}}
\]
或
\phantomsection\label{eq:9-6}
\begin{equation}
h_i^2=l_{i1}^2+l_{i2}^2+\cdots+l_{im}^2
\tag{9-6}
\end{equation}
和
\[
\sigma_{ii}=h_i^2+\psi_i,\qquad i=1,2,\cdots,p
\]
第 $i$ 个共性方差是第 $i$ 变量在这 $m$ 个公共因子上载荷的平方和.

\noindent\textbf{例 9.1}\quad (对两因子验证关系 $\boldsymbol{\Sigma}=\mathbf{L}\mathbf{L}'+\boldsymbol{\Psi}$)

考虑协方差阵
\[
\boldsymbol{\Sigma}=\begin{bmatrix}
19&30&2&12\\
30&57&5&23\\
2&5&38&47\\
12&23&47&68
\end{bmatrix}
\]
等式
\[
\begin{bmatrix}
19&30&2&12\\
30&57&5&23\\
2&5&38&47\\
12&23&47&68
\end{bmatrix}
=\begin{bmatrix}4&1\\7&2\\-1&6\\1&8\end{bmatrix}
\begin{bmatrix}4&7&-1&1\\1&2&6&8\end{bmatrix}
+\begin{bmatrix}
2&0&0&0\\
0&4&0&0\\
0&0&1&0\\
0&0&0&3
\end{bmatrix}
\]
